% ===========================================================================
% A more powerful selective inference test for k-means clustering
% Skeleton draft -- structure mirrors Chen, Jewell & Witten (2023),
% "More powerful selective inference for the graph fused lasso" (JCGS).
%
% Every empirical number is pulled from the repository:
%   - memory/union-method-status.md, memory/figure-system.md
%   - sims/results/datathin_compare.rds, power_type1_summary.rds,
%     repro_yunbarber_q{2,10}.rds, penguins_pvalues_k3.csv, scrna_pvalues_k3.csv
% Anything not yet verified is flagged with a "% TODO:" comment.
% Do NOT treat un-TODO'd numbers as placeholders -- they are real.
% ===========================================================================
\documentclass[11pt]{article}

\usepackage{amsmath,amssymb,amsthm}
\usepackage{natbib}
\usepackage{graphicx}
\usepackage{booktabs}
\usepackage[margin=1in]{geometry}
\usepackage{bm}
\usepackage{xcolor}
\usepackage{hyperref}

\graphicspath{{../sims/results/}}

\newtheorem{theorem}{Theorem}
\newtheorem{proposition}{Proposition}
\newtheorem{lemma}{Lemma}
\newtheorem{definition}{Definition}

\newcommand{\R}{\mathbb{R}}
\newcommand{\E}{\mathbb{E}}
\newcommand{\Prob}{\mathbb{P}}
\newcommand{\X}{\mathbf{X}}
\newcommand{\x}{\mathbf{x}}
\newcommand{\bmu}{\bm{\mu}}
\newcommand{\bnu}{\bm{\nu}}
\newcommand{\Cset}{\mathcal{C}}
\newcommand{\Sset}{\mathcal{S}}
\newcommand{\Pp}{\mathbf{P}}
\newcommand{\norm}[1]{\left\lVert #1 \right\rVert}
\newcommand{\Fnorm}[1]{\left\lVert #1 \right\rVert_F}
\newcommand{\todo}[1]{\textcolor{red}{[\textbf{TODO:} #1]}}

\title{A More Powerful Selective Inference Test for \\ $k$-Means Clustering}

\author{%
  % TODO: author list and affiliations
  Yiqun Chen \\ \textit{Draft skeleton --- \today}
}

\date{}

\begin{document}
\maketitle

\begin{abstract}
After partitioning data into clusters with $k$-means, a natural next step is to
test whether two estimated clusters have different means. Because the very
hypothesis is chosen by looking at the data, a classical $z$- or $\chi^2$-test is
badly anti-conservative: on the Palmer penguins data its $p$-values for three
cluster pairs are on the order of $10^{-57}$, $10^{-29}$, and $10^{-134}$.
\citet{chen2023kmeans} solve this by conditioning on the \emph{entire} realized
sequence of Lloyd-algorithm cluster assignments and computing an exact truncated
$\chi$ $p$-value. Conditioning on the full algorithmic path is convenient but
discards power. We make three contributions. First, following the
``condition on less'' principle of \citet{chen2023gfl} and the selective
optimality theory of \citet{fithian2017optimal}, we propose a test that
conditions only on the event that the two \emph{tested} clusters are preserved,
taken as a union over all Lloyd paths that produce them; this enlarges the
truncation set and is uniformly more powerful, with the gain widening as the
number of clusters $K$ grows. Second, when the noise variance is unknown, we
develop an exact studentized test built on a Fisher--$F$ pivot, generalizing
\citet{yun2023unknownvar} from $K=2$ to arbitrary $K$ and eliminating their
Monte Carlo step; the test requires no estimate of $\sigma$ at all. Third, we
give a fast numerical/analytic solver for the quadratic truncation inequalities
that define the conditioning set
\citep{yun2024multiplepairs}. In simulations the union test controls the
selective Type I error (rejection rate $0.035$--$0.045$ at the global null versus
$0.975$ for the naive test) while substantially improving power, and a head-to-head
comparison against data thinning \citep{neufeld2024datathin} shows that our
studentized test detects real clusters far more often --- because data thinning
must spend half the data on the test fold and must estimate $\sigma$ to perform
the split, whereas our test needs neither.
\end{abstract}

% ===========================================================================
\section{Introduction}
% ===========================================================================

Clustering is routinely used to generate hypotheses: one partitions the
observations, notices that two groups look different, and then asks whether the
difference is real. The trouble is that the groups, and hence the hypothesis,
were chosen using the same data used to test it --- a form of double dipping.
A classical Wald test that ignores this selection has a wildly inflated Type I
error rate \citep{gao2024hierarchical,chen2023kmeans}. The size of the problem is
easy to see empirically: clustering the Palmer penguins data into $K=3$ groups
and applying a naive $\chi^2$ test to the three cluster pairs gives $p$-values of
$3.3\times10^{-57}$, $7.5\times10^{-29}$, and $5.6\times10^{-134}$ (Section~\ref{sec:realdata});
these are absurd certificates of significance produced \emph{purely} by the act
of clustering.

Selective inference repairs this by conditioning on the selection event
\citep{lee2016exact,fithian2017optimal}. For $k$-means clustering,
\citet{chen2023kmeans} test
$H_0: \bmu^\top\bnu = \bm 0$ for the contrast $\bnu$ encoding the difference in
means of two estimated clusters, conditioning on the clustering output (and on
nuisance directions) so that the conditional law of the test statistic is a
truncated $\chi$ distribution. Their test controls the selective Type I error and
is computable in closed form. However, it conditions on the \emph{full realized
Lloyd path}: the entire sequence of intermediate cluster assignments produced by
the algorithm, iteration by iteration. This is far more information than the
hypothesis requires --- the hypothesis only concerns the two clusters that were
ultimately tested --- and conditioning on it shrinks the truncation set and the
power.

\paragraph{Contributions.} This paper makes three contributions.
\begin{enumerate}
  \item \textbf{A more powerful test that conditions on less}
        (Section~\ref{sec:union}). We condition only on the event that the two
        \emph{tested} clusters are preserved, unioned over every Lloyd path that
        outputs them. This is the $k$-means analogue of the
        ``condition on less $\Rightarrow$ more power'' principle that
        \citet{chen2023gfl} used for the graph fused lasso, and is justified by
        the selective-optimality framework of \citet{fithian2017optimal}. The
        resulting truncation set is a superset of the path set, so the test is
        uniformly more powerful; empirically the power gain over the path test
        widens with the number of clusters $K$.
  \item \textbf{An exact studentized test for unknown variance}
        (Section~\ref{sec:rfiber}). Following \citet{yun2023unknownvar} we
        studentize, obtaining an $F$ pivot in which $\sigma$ cancels. We
        characterize the conditioning set \emph{exactly} on the
        ``$R$-fiber'' --- a one-parameter rotation of the data at fixed total
        energy --- for any $K$, replacing the importance sampling that
        \citet{yun2023unknownvar} require for $K>2$ with an exact, zero-variance
        $p$-value.
  \item \textbf{A fast solver} (Section~\ref{sec:computation}). The conditioning
        set is the solution set of a collection of quadratic inequalities in the
        perturbation parameter; we compute the coefficients analytically and
        solve by an interval/scan method, in the spirit of
        \citet{yun2024multiplepairs}.
\end{enumerate}

We also report a head-to-head comparison against data thinning
\citep{neufeld2024datathin,dharamshi2024generalized}, a sample-splitting
alternative to selective inference. Data thinning is valid, but it pays two
prices that our studentized test avoids: it clusters on only half the data, so it
detects fewer real clusters, and it needs an estimate of $\sigma$ to perform the
split --- an estimate that is inflated two- to three-fold under genuine cluster
separation, poisoning the split.

% ===========================================================================
\section{Background: the path test of Chen and Witten}
\label{sec:background}
% ===========================================================================

\subsection{Model and hypothesis}
We observe $\X\in\R^{n\times q}$ with independent rows,
\begin{equation}
  \X \sim \mathcal{MN}_{n\times q}(\bmu,\, \mathbf I_n,\, \sigma^2\mathbf I_q),
  \qquad\text{i.e.}\quad X_i \sim \mathcal N(\bmu_i,\sigma^2\mathbf I_q),
\end{equation}
with unknown mean rows $\bmu_i$ and (for now) known $\sigma^2$.
Running $k$-means with $K$ clusters returns a partition
$\hat\Cset_1,\dots,\hat\Cset_K$ of $\{1,\dots,n\}$. Given two estimated clusters
$\hat\Cset_1,\hat\Cset_2$ we test
\begin{equation}
  H_0:\ \tfrac{1}{|\hat\Cset_1|}\!\sum_{i\in\hat\Cset_1}\!\bmu_i
       = \tfrac{1}{|\hat\Cset_2|}\!\sum_{i\in\hat\Cset_2}\!\bmu_i
  \quad\Longleftrightarrow\quad
  H_0:\ \bmu^\top\bnu = \bm 0_q,
\end{equation}
where the contrast $\bnu\in\R^n$ has entries
$\nu_i = \mathbf 1\{i\in\hat\Cset_1\}/|\hat\Cset_1| - \mathbf 1\{i\in\hat\Cset_2\}/|\hat\Cset_2|$.
The natural statistic is $\norm{\X^\top\bnu}_2$, the length of the estimated
difference in means.

\subsection{The selective $p$-value and the Lloyd path}
The naive Wald $p$-value
$p_{\mathrm{naive}} = \Prob(\norm{\X^\top\bnu}_2 \ge \norm{\x^\top\bnu}_2)$
ignores that $\bnu$ was chosen by clustering and is severely anti-conservative.
\citet{chen2023kmeans} instead condition on the clustering output. Writing
$c_i^{(t)}(\X)$ for the cluster assigned to observation $i$ at iteration $t$ of
Lloyd's algorithm, they define
\begin{equation}
  p_{\mathrm{path}} = \Prob_{H_0}\!\left(
    \norm{\X^\top\bnu}_2 \ge \norm{\x^\top\bnu}_2 \ \Big|\
    \bigcap_{t=0}^{T}\bigcap_{i=1}^{n}\!\big\{c_i^{(t)}(\X)=c_i^{(t)}(\x)\big\},\
    \Pp_\nu^\perp\X = \Pp_\nu^\perp\x,\
    \mathrm{dir}(\X^\top\bnu)=\mathrm{dir}(\x^\top\bnu)
  \right).
\end{equation}
The conditioning event fixes everything except a single scalar: perturbing the
data along $\bnu\,\mathrm{dir}(\x^\top\bnu)^\top$,
$\x'(\phi) = \x + \big((\phi-\norm{\x^\top\bnu}_2)/\norm\bnu_2^2\big)\,\bnu\,\mathrm{dir}(\x^\top\bnu)^\top$,
the $p$-value reduces to a one-dimensional truncated-$\chi$ probability,
\begin{equation}
  p_{\mathrm{path}} = \Prob\!\left(\phi \ge \norm{\x^\top\bnu}_2 \mid \phi\in\Sset_T\right),
  \qquad \phi \sim (\sigma\norm\bnu_2)\cdot\chi_q,
\end{equation}
where $\Sset_T = \{\phi\ge 0 : \text{the entire Lloyd path is preserved at }\x'(\phi)\}$.
The key computational fact is that $\Sset_T$ is the intersection of the solution
sets of $O(nKT)$ quadratic inequalities in $\phi$, so it is a union of intervals
computable analytically; the test controls the selective Type I error in the
sense of \citet{lee2016exact} and \citet{fithian2017optimal}.

\subsection{The orthogonal decomposition}
\label{sec:decomp}
A geometric view, which we use throughout, decomposes the data with three
Frobenius-orthogonal projections in $\R^{n\times n}$,
\begin{equation}
  \X = \Pp_0\X + \Pp_1\X + \Pp_2\X,
\end{equation}
where (i) the rank-one \emph{between-cluster} projection $\Pp_0=\bnu\bnu^\top/\norm\bnu_2^2$
carries the contrast, with $\Fnorm{\Pp_0\X}=\norm{\X^\top\bnu}_2/\norm\bnu_2$;
(ii) the \emph{within-cluster} centering projection $\Pp_1$ of the two tested
clusters (rank $m-2$, with $m=|\hat\Cset_1|+|\hat\Cset_2|$) carries the
within-cluster scatter; and (iii) the \emph{nuisance} $\Pp_2=\mathbf I-\Pp_0-\Pp_1$.
The path test moves only $\Fnorm{\Pp_0\X}$ along a ray, holding $\Pp_1\X$ and
$\Pp_2\X$ fixed; Figure~\ref{fig:decomposition} contrasts this with the
studentized test of Section~\ref{sec:rfiber}.

\textbf{Why the path test loses power.} Conditioning on the whole path is
\emph{sufficient} for validity but not \emph{necessary}: the hypothesis depends
only on the two tested clusters, not on the order in which Lloyd's algorithm
reached them. Many distinct intermediate assignment sequences end in the same
final partition of the two tested clusters; the path test discards all of them
but the realized one. By \citet{fithian2017optimal}, conditioning on this extra,
ancillary information can only cost power.

% ===========================================================================
\section{A more powerful test: condition on less}
\label{sec:union}
% ===========================================================================

\subsection{The union conditioning set}
Our test conditions on the \emph{final-partition} event for the two tested
clusters, rather than the full path. Define
\begin{equation}
  \Sset_{\mathrm{union}}
  = \big\{\phi\ge 0:\ \text{$k$-means on $\x'(\phi)$ outputs clusters that
    preserve both $\hat\Cset_1$ and $\hat\Cset_2$}\big\}.
\end{equation}
Concretely, along the perturbation $\x'(\phi)$ we re-run $k$-means (with a shared
initialization) and record, for each $\phi$, whether the two tested clusters
reappear --- regardless of which Lloyd path produced them. Because the realized
path is one of the paths that preserve the tested clusters,
\begin{equation}
  \Sset_T \ \subseteq\ \Sset_{\mathrm{union}},
\end{equation}
so $\Sset_{\mathrm{union}}$ is a superset of the path set, and the corresponding
truncated-$\chi$ $p$-value,
$p_{\mathrm{union}} = \Prob(\phi \ge \norm{\x^\top\bnu}_2 \mid \phi\in\Sset_{\mathrm{union}})$,
is based on weaker conditioning. This is exactly the construction
\citet{chen2023gfl} use for the graph fused lasso, transported to $k$-means:
condition on the \emph{tested object} (here, the two clusters) instead of the
\emph{algorithmic trace} (the Lloyd path).

\subsection{Validity}
The union set is, like the path set, a function of $\phi$ alone along the
perturbation line (everything orthogonal to the contrast is held fixed), so
$\phi \mid \phi\in\Sset_{\mathrm{union}}$ is again a truncated scaled-$\chi_q$
and the test controls the selective Type I error.
% TODO: state and prove this as a formal proposition (the conditioning event is
% a deterministic function of phi; truncated-chi remains a valid pivot under any
% phi-measurable conditioning set). Mirror Prop 1 of Chen, Jewell & Witten (2023).
Empirically, at the global null the union test rejects at rate $0.045$
(KS-uniform), versus $0.975$ for the naive test
(Table~\ref{tab:typeI}; \texttt{power\_type1\_summary.rds}).

\subsection{Construction of the union set}
We enumerate the distinct Lloyd paths reachable along the perturbation line as
follows. Starting from the observed $\phi=\norm{\x^\top\bnu}_2$, we compute the
$\phi$-interval on which the observed path is realized (the path test's
$\Sset_T$), then probe just past each interval endpoint, re-run $k$-means there
to discover the neighbouring path, compute \emph{its} preservation interval, and
repeat, queueing newly discovered endpoints. The union of all intervals whose
path preserves both tested clusters is $\Sset_{\mathrm{union}}$; the subset of
intervals matching the observed path signature recovers $\Sset_T$ exactly, so the
two $p$-values are produced from a single sweep. The exploration is purely
Euclidean (it perturbs $\x$ and re-runs Euclidean $k$-means), which makes it
\emph{metric-independent}: the known general-covariance test
$\X\sim\mathcal{MN}(\bmu,\mathbf I_n,\bm\Sigma)$ reuses the identical exploration
and only rescales the resulting intervals by a scalar and swaps the test
statistic for $\sqrt{(\x^\top\bnu)^\top\bm\Sigma^{-1}(\x^\top\bnu)}$. The
general-covariance path $p$-value matches the isotropic one to machine precision.

\paragraph{A scan cutoff.} The far tail of the perturbation line contributes
negligible $\chi$ mass to the $p$-value. We therefore stop exploring once $\phi$
exceeds the point where the conditional $\chi$ survival mass beyond it falls below
a fraction \texttt{scan\_tail\_prob} (default $10^{-6}$) of the survival mass
beyond the observed statistic. This leaves the $p$-value essentially unchanged
(verified $\Delta \le 10^{-11}$) while cutting roughly $25$--$40\%$ of the path
exploration.

% ===========================================================================
\section{Unknown variance: the studentized $R$-fiber}
\label{sec:rfiber}
% ===========================================================================

\subsection{The $F$ pivot}
When $\sigma$ is unknown we cannot compare $\norm{\x^\top\bnu}_2$ to a scaled
$\chi_q$, because the scale is exactly what we lack. Following
\citet{yun2023unknownvar} we studentize. Using the decomposition of
Section~\ref{sec:decomp}, the within-cluster scatter $\Fnorm{\Pp_1\X}^2$ is an
unbiased $\sigma^2$-estimate independent of the between-cluster signal
$\Fnorm{\Pp_0\X}^2$, so their ratio is pivotal:
\begin{equation}
  R \;=\; \frac{\Fnorm{\Pp_0\X}^2/q}{\Fnorm{\Pp_1\X}^2/\big((m-2)q\big)}
     \;=\; (m-2)\,\frac{\Fnorm{\Pp_0\X}^2}{\Fnorm{\Pp_1\X}^2}
     \;\sim\; F_{q,\,(m-2)q} \quad\text{under } H_0',
  \label{eq:Rstat}
\end{equation}
where $m=|\hat\Cset_1|+|\hat\Cset_2|$ counts only the two tested clusters and the
unknown $\sigma$ cancels. The required null is the slightly stronger
\emph{pointwise} null $H_0'$: within each tested cluster the means $\bmu_i$ are
equal to the cluster mean (so that $\Fnorm{\Pp_1\X}^2$ is a clean noise estimate).
This is the price of not knowing $\sigma$.
% TODO: discuss H0' vs H0 explicitly and note the adversarial finding: broad-null-
% but-not-pointwise violations (equal cluster means, canceling +/- sub-populations)
% only inflate the denominator ||P1X||, so the test is CONSERVATIVE; anti-
% conservatism requires within-cluster heterogeneity to deflate ||P1X|| under a
% TRUE mean difference, which k-means tends to resolve by spending a cluster.
% (source: union-method-status.md)

\subsection{The rotation and the fiber}
We condition on the clustering, on $\Pp_2\X$, on the unit directions
$U_0=\Pp_0\X/\Fnorm{\Pp_0\X}$ and $U_1=\Pp_1\X/\Fnorm{\Pp_1\X}$ (which are
orthonormal), and on the \emph{total energy}
$T = \Fnorm{\Pp_0\X}^2+\Fnorm{\Pp_1\X}^2$. The single remaining degree of freedom
is an angle $\theta$, and the conditional perturbation is
\begin{equation}
  \x(\theta) \;=\; \Pp_2\x \;+\; \sqrt{T}\,\big(\sin\theta\,U_0 + \cos\theta\,U_1\big),
  \qquad \x(\theta_{\mathrm{obs}})=\x,
  \label{eq:fiber}
\end{equation}
with $\theta_{\mathrm{obs}}=\operatorname{atan2}(\Fnorm{\Pp_0\x},\Fnorm{\Pp_1\x})$
and $R(\theta)=(m-2)\tan^2\theta$. Because $U_0\perp U_1$ are orthonormal,
\eqref{eq:fiber} is a \textbf{Givens rotation} of the vectorized data
$\x\in\R^{nq}$ inside the plane $\mathrm{span}\{U_0,U_1\}$, fixing its orthogonal
complement,
\begin{equation}
  \x(\theta) = G(\theta-\theta_{\mathrm{obs}})\,\x,\qquad
  G(\delta) = \mathbf I + (\cos\delta-1)\big(U_0U_0^\top+U_1U_1^\top\big)
                       + \sin\delta\,\big(U_0U_1^\top-U_1U_0^\top\big).
\end{equation}
Since $G$ is orthogonal it preserves the in-plane radius, so $T$ is held fixed and
only the ratio $R$ varies --- which is precisely why the test is free of $\sigma$:
$T$ is complete and sufficient for $\sigma$ under $H_0'$ (Basu's theorem), so the
scale-invariant ratio $R$ is ancillary to it. Contrast this with the known-$\sigma$
ray $\x'(\phi)$, which moves $\Fnorm{\Pp_0\X}$ in a straight line while
$\Pp_1\X$ stays put. The two slices are drawn in Figure~\ref{fig:decomposition}:
the known-$\sigma$ test walks a horizontal ray (fix within-scatter, vary
separation $\to\chi$), the studentized test walks an arc (fix total energy, vary
the angle $\to F$).

\paragraph{Geometry.} Datasets of fixed energy $T$ lie on a sphere of radius
$\sqrt T$. The pole ($\theta=\pi/2$) is all separation (within-scatter $0$), the
equator ($\theta=0$) is fully merged clusters (no separation), the latitude is
$\theta$ (and $R=(m-2)\tan^2\theta$ depends only on latitude), and the longitude
indexes the within-scatter pattern. Our test fixes the longitude and sweeps
$\theta$: it walks a single meridian, the latitude playing the role of the test
statistic.

\subsection{The truncated-$F$ test and the union}
With $\Sset = \{\theta\in(0,\pi/2):\text{the clustering is preserved at }\x(\theta)\}$,
the selective $p$-value is the truncated-$F$ probability
\begin{equation}
  p = \Prob\big(R(\theta)\ge R \mid \theta\in\Sset\big),
  \qquad R\sim F_{q,(m-2)q}.
\end{equation}
The cluster-assignment margins $\norm{x_i-m_j}^2-\norm{x_i-m_k}^2$ along the
fiber are quadratic forms in $(\sin\theta,\cos\theta)$, so $\Sset$ is a union of
arcs computed exactly (Section~\ref{sec:computation}). Unioning the arcs across
all Lloyd paths that preserve the two tested clusters yields the more powerful
union version --- and, crucially, this union is computed \emph{on the same
$R$-fiber}, so the truncated-$F$ remains a valid pivot over it.

\paragraph{An important caveat: do not remap the known-$\sigma$ union.}
A natural shortcut is to take the known-$\sigma$ $\phi$-union set and feed it,
re-scaled, into the $F$ pivot. This is \emph{invalid}: the studentized-$F$ pivot
holds only \emph{per Lloyd path}, and the per-path truncated-$F$'s do not
aggregate into a union-$F$ unless the union is taken in the $F$ pivot's own
coordinate. Along the $\phi$-ray the within-scatter $\Fnorm{\Pp_1\X}$ is held
fixed, but along the $R$-fiber it varies, so the per-path cancellation that makes
$R$ pivotal breaks when one imports the $\phi$-fiber's set. Exhaustive
simulation confirms the remapped union over-rejects (rejection $0.081$,
scaling with union width from $0.010$ up to $0.170$), whereas the same union taken
on the $R$-fiber is calibrated. The valid construction perturbs along the
$R$-fiber \eqref{eq:fiber} and unions the $\theta$-arcs, so that
$\{\text{clusters preserved}\} = \{R\in\Sset'\}$ is a function of $R$ alone and
the union of a one-dimensional truncation is exact.
% TODO: this caveat is methodologically central -- decide how much to foreground
% it. It is the single most error-prone part of the method.

\paragraph{Generalizing Yun and Barber to any $K$.}
For $K=2$ the two tested clusters are all the data ($m=n$) and
\citet{yun2023unknownvar} compute the truncation set exactly by reusing
\citet{chen2023kmeans}. For $K>2$ they resort to importance sampling. Our exact
$R$-fiber solver computes the arc set exactly for any $K$, so it
(i) replaces their Monte Carlo with an exact, zero-variance $p$-value for $K>2$
$k$-means (even for the path test), and (ii) yields the more powerful union for
free. The method is $k$-means-specific (theirs applies to general clustering,
e.g. hierarchical) and requires the pointwise null $H_0'$.

\paragraph{When does the union help under unknown variance?}
The union's power gain is largely a \emph{known-variance} phenomenon. A one-sided
truncated $p$-value shrinks when enlarging the truncation set only if the added
mass lands \emph{below} the observed statistic. For the known-$\sigma$ $\chi$
statistic, alternative Lloyd paths tolerate smaller separation and so extend the
$\phi$-interval downward, below the observed value --- producing large gains. For
the studentized $F$, $k$-means has \emph{minimized} the within-scatter
$\Fnorm{\Pp_1\X}$, pinning $R_{\mathrm{obs}}$ against its low-$R$ floor, so the
extra room is almost entirely \emph{above} the observed statistic and the union
$p$-value $\approx$ the path $p$-value. In short, the union's leverage is
\emph{scale} (contrast-length) information, which a scale-free statistic discards.
This is why we report the union gain primarily under known $\sigma$
(Figure~\ref{fig:cwpower}) and report the studentized test primarily as the
$\sigma$-free route to validity and to the Yun--Barber power regime
(Figure~\ref{fig:yunbarber}).

% ===========================================================================
\section{Computation}
\label{sec:computation}
% ===========================================================================

Both conditioning sets reduce to solving quadratic inequalities in a single
perturbation parameter.

\paragraph{Known $\sigma$ ($\phi$-ray).} Along $\x'(\phi)$ each $k$-means
assignment margin is a quadratic in $\phi$, so the per-path preservation set is
the intersection of $O(nKT)$ quadratic inequalities, computed in closed form, and
the union set is assembled by the endpoint-probing sweep of
Section~\ref{sec:union}.

\paragraph{Unknown $\sigma$ ($R$-fiber).} Along $\x(\theta)$ each assignment
margin is an exact quadratic form
$Q(a,b)=Aa^2+Bb^2+Cab+Da+Eb+F_c$ in $(a,b)=(\sin\theta,\cos\theta)$.
Substituting $t=\tan(\theta/2)$ turns $Q(\sin\theta,\cos\theta)(1+t^2)^2$ into a
\emph{quartic} $N(t)$, whose real roots in $[0,1]$ give the arc breakpoints; the
sign of $Q$ at a midpoint decides which arcs satisfy the inequality, and
intersecting over all margins yields the arc set for one Lloyd trajectory
(\texttt{sims/rfiber\_exact.R}). The truncated-$F$ mass over the resulting
$r$-intervals is computed in a cancellation-proof way (taking the maximum of the
CDF-difference and the survival-difference, the latter stable in the extreme
upper tail under strong separation). The solver is modular and unit-tested:
$23$ tests for the $F$-region routine (including agreement with Monte Carlo and
extreme-tail underflow), and the arc assembler passes its trajectory-preservation
unit tests; the truncated-$F$ pivot is verified uniform under fixed truncation.
% TODO: confirm the exact unit-test counts before publication. union-method-status.md
% records 14/14 for the 7 primitive modules, 18/18 for the path_arc assembler,
% and figure-system.md records 23/23 for the F-region solver (sims/test_f_region.R).
A pure interval/scan implementation also computes the set without analytic
quartics, agreeing with the exact solver where the grid is fine and finding more
arcs where the grid is coarse; numerically these two solvers agree.
The same numerical-solver idea underlies the multiple-pairs extension of
\citet{yun2024multiplepairs}.

% ===========================================================================
\section{Simulations}
\label{sec:sims}
% ===========================================================================

Unless noted, data are $q$-dimensional Gaussian with isotropic unit variance;
$k$-means uses a fixed shared initialization. We report rejection rates at level
$0.05$.

\subsection{Type I error}
At the global null ($\bmu=\bm 0$) the union, path, and studentized tests are all
calibrated, while the naive test is invalid. From the global-null cell of
\texttt{power\_type1\_summary.rds} ($200$ replicates), the union rejects at
$0.045$ and the path at $0.060$; the naive test rejects at $0.975$
(Table~\ref{tab:typeI}). The studentized $R$-fiber test is likewise calibrated
across $q\in\{2,10,50,100\}$ (Type I $0.047/0.046/0.029/0.054$; slightly
conservative at high $q$). Under AR$(0.5)$/AR$(0.9)$ covariance the
general-covariance union and path remain on the diagonal while the naive test
flatlines (Figure~\ref{fig:gencov}). Under heavy-tailed $t_5$/$t_{10}$ noise the
studentized test is fairly robust, with mild anti-conservatism only for the
heaviest tails (Figure~\ref{fig:heavytail}).

\begin{table}[t]
\centering
\caption{Selective Type I error at the global null ($\bmu=\bm 0$).
Rejection rates at level $0.05$. The naive test is grossly invalid; the
selective tests are calibrated. Known-$\sigma$ union/path from
\texttt{power\_type1\_summary.rds} ($n=20K$, $q=2$, $200$ reps); studentized
($R$-fiber) and data-thinning columns from the null cell of
\texttt{datathin\_compare.rds} ($K=3$, $q=2$, $400$ reps).}
\label{tab:typeI}
\begin{tabular}{lccccc}
\toprule
 & naive & path (known) & union (known) & studentized $F$ & data thinning (oracle) \\
\midrule
rejection rate & 0.975 & 0.045--0.060 & 0.035--0.045 & 0.040 & 0.065 \\
\bottomrule
\end{tabular}
\end{table}

\subsection{Power: union versus path (known $\sigma$)}
\label{sec:powersim}
Table~\ref{tab:power} reports power for the known-$\sigma$ union and path tests as
a function of the separation $\delta$ (from \texttt{power\_type1\_summary.rds}).
The union is uniformly more powerful, with gains of order $0.1$--$0.3$; averaged
over $\delta>0$ the union improves on the path by about $16.8$ percentage points.
The gain \emph{widens with $K$}: at $\delta=5$, $q=2$ the union's advantage over
the path is $0.075/0.115/0.225/0.200$ for $K=2/3/4/5$. This is a distinctive
feature of our exact solver, which handles any $K$.
% TODO: the K=2..5 gains come from figure-system.md (the dropped varying_k figure,
% now folded into standard_sweep). Confirm against sims/results/standard_sweep data
% before final.

\begin{table}[t]
\centering
\caption{Power of the known-$\sigma$ union vs.\ path test against separation
$\delta$ ($n=20K$, $q=2$, $K=3$; \texttt{power\_type1\_summary.rds}, $200$ reps).
The union is uniformly more powerful. Figure~\ref{fig:cwpower}.}
\label{tab:power}
\begin{tabular}{lccccccc}
\toprule
$\delta$ & 0 & 1 & 2 & 3 & 4 & 5 & 6 \\
\midrule
union & 0.045 & 0.075 & 0.195 & 0.445 & 0.660 & 0.910 & 0.925 \\
path  & 0.060 & 0.060 & 0.085 & 0.230 & 0.415 & 0.645 & 0.770 \\
\bottomrule
\end{tabular}
\end{table}

\subsection{Unknown variance: reproducing the Yun--Barber regime}
We reproduce the equilateral-triangle, $K=3$, $q\in\{2,10\}$ design of
\citet{yun2023unknownvar} in our $k$-means framework (path test, clusters $1$ vs
$2$, $300$ reps; \texttt{repro\_yunbarber\_q\{2,10\}.rds}), comparing four
variance handlings on the \emph{same} truncation set: an oracle (known $\sigma$),
the studentized $F$, a median-based $\hat\sigma_{\mathrm{MED}}$ plug-in, and a
sample-SD plug-in. The ordering oracle $\ge$ studentized $\ge$ plug-ins holds
throughout. The studentized $F$ decisively beats the sample plug-in once the
clusters separate: at $q=2$, $\delta=7$, power is $0.490$ (studentized) versus
$0.290$ (sample) and $0.717$ (oracle); the gap opens at $\delta\approx 4$--$5$.
The mechanism is that the global sample SD absorbs the between-cluster
separation, so $\E[\hat\sigma_{\mathrm{sample}}]$ climbs roughly linearly in
$\delta$ and the plug-in power plateaus, whereas the within-cluster scale stays
near $1$ and the $F$ approaches the oracle. The robust
$\hat\sigma_{\mathrm{MED}}$ is dimension-dependent: near-oracle at $q=10$ but, with
three balanced clusters, as inflated as the sample SD at $q=2$.

\subsection{Comparison with data thinning}
\label{sec:datathin}
Data thinning \citep{neufeld2024datathin} is a sample-splitting alternative:
split each entry $X=X^{(1)}+X^{(2)}$ into independent Gaussian halves, cluster on
$X^{(1)}$, and test on $X^{(2)}$ with an ordinary (unconditional) mean test. It is
valid by construction, but it pays two prices our studentized test avoids.
\emph{First}, it clusters on only half the information, so it recovers the true
pair less often. \emph{Second}, to perform the split it needs a variance: the
oracle ($\sigma^2$) is unavailable in practice, and any plug-in is inflated under
cluster separation, which poisons both folds.

Table~\ref{tab:datathin} reports the head-to-head, per-replicate matched, on a
$K=3$ equilateral-triangle design ($q=2$, $n_{\mathrm{per}}=10$,
$\varepsilon=0.5$; \texttt{datathin\_compare.rds}, null $400$ reps, signal $220$
reps each). Three things stand out. (i) At the global null all selective tests
and oracle thinning control Type I; the plug-in thinning variants are slightly
anti-conservative ($0.075$ for the MED split). (ii) \emph{Detection probability}
--- whether the clustering recovers the true pair --- is far higher for the
full-data selective pipeline (e.g.\ $0.91$ at $\delta=4$) than for plug-in
thinning, which clusters on half the data \emph{and} with an inflated split
variance, collapsing detection to $0.08$--$0.23$. (iii) Even \emph{conditional} on
detection, the selective union and studentized tests are competitive with or
better than plug-in thinning, while needing no $\sigma$. The recorded variance
estimates make the poisoning explicit: under separation the MED and sample
estimates inflate to $2.08$--$3.12$ (true $\sigma=1$) as $\delta$ runs $4\to 7$.
Only oracle thinning --- which knows the true $\sigma$ and so is not realizable
--- keeps high detection. Our studentized test needs no $\sigma$ at all and
clusters on the full data. Figure~\ref{fig:datathin} draws all three metrics in
the paper's shared grammar (colour $=$ variance handling; line style $=$
paradigm, selective solid vs.\ thinning dashed): panel~A shows the Type I
calibration as a QQ plot against the uniform, where every valid method ---
selective \emph{and} thinning --- tracks the $45^\circ$ line while the naive test
plunges below it, and panels~B--C plot detection probability and conditional
power against $\delta$.

\begin{figure}[t]\centering
  \includegraphics[width=\textwidth]{datathin.png}
  \caption{Data thinning vs.\ the selective tests, all three metrics in the
  paper's shared grammar (colour $=$ variance handling; line style $=$ paradigm,
  selective solid vs.\ data thinning dashed). \textbf{A:} Type I calibration as a
  QQ plot against the uniform --- every valid method, selective \emph{and}
  thinning, tracks the $45^\circ$ line, while the naive test (orange) plunges
  below it. \textbf{B:} detection probability vs.\ separation $\delta$ ---
  full-data selective clustering (solid black) dominates the thinned splits, and
  the plug-in splits ($\hat\sigma_{\mathrm{MED}}$ green, $\hat\sigma_{\mathrm{sample}}$
  vermilion) collapse because the variance estimate is inflated by the very
  separation being tested. \textbf{C:} conditional power vs.\ $\delta$ (reject
  $\mid$ detected); oracle thinning rejects whenever it detects, but rarely
  detects. Cf.\ Table~\ref{tab:datathin}. Source:
  \texttt{sims/plot\_datathin.R} $\to$ \texttt{datathin.png}.}
  \label{fig:datathin}
\end{figure}

\begin{table}[t]
\centering
\caption{Data thinning vs.\ the selective tests, matched per replicate
(\texttt{datathin\_compare.rds}; $K=3$, $q=2$, $n_{\mathrm{per}}=10$,
$\varepsilon=0.5$). ``detect'' is the probability the true pair is recovered;
``cpow'' is rejection rate \emph{conditional} on detection. Oracle thinning uses
the true $\sigma$ (not realizable); MED/sample thinning estimate $\sigma$ to split.
The median split-variance estimates ($\hat\sigma_{\mathrm{MED}}$, true $\sigma=1$)
show the inflation that poisons the folds.}
\label{tab:datathin}
\begin{tabular}{lcccc}
\toprule
$\delta$ & 4 & 5 & 6 & 7 \\
\midrule
$\hat\sigma_{\mathrm{MED}}$ (true $=1$)        & 2.08 & 2.42 & 2.77 & 3.12 \\
\midrule
\textbf{detect} selective (full data)          & 0.91 & 0.94 & 0.85 & 0.84 \\
\textbf{detect} thinning, oracle $\sigma$      & 0.60 & 0.88 & 0.93 & 0.95 \\
\textbf{detect} thinning, MED split            & 0.08 & 0.12 & 0.18 & 0.23 \\
\textbf{detect} thinning, sample split         & 0.11 & 0.11 & 0.15 & 0.17 \\
\midrule
\textbf{cpow} union (known $\sigma$)           & 0.79 & 0.92 & 0.99 & 0.99 \\
\textbf{cpow} studentized $F$ (no $\sigma$)    & 0.37 & 0.62 & 0.87 & 0.96 \\
\textbf{cpow} thinning, oracle $\sigma$        & 1.00 & 1.00 & 1.00 & 1.00 \\
\textbf{cpow} thinning, MED split              & 0.44 & 0.63 & 0.65 & 0.60 \\
\textbf{cpow} thinning, sample split           & 0.50 & 0.32 & 0.56 & 0.51 \\
\bottomrule
\end{tabular}
\end{table}

% ===========================================================================
\section{Real data}
\label{sec:realdata}
% ===========================================================================

\subsection{Palmer penguins}
We cluster the Palmer penguins data \citep{horst2020palmerpenguins} --- four
standardized morphological features, the same data used by
\citet{yun2023unknownvar} --- into $K=3$ groups and test all three cluster pairs
(\texttt{penguins\_pvalues\_k3.csv}, $\hat\sigma_{\mathrm{MED}}=1.155$).
Table~\ref{tab:penguins} shows the pattern. The naive test reports absurd
$p$-values ($3.3\times10^{-57}$, $7.5\times10^{-29}$, $5.6\times10^{-134}$). The
original known-$\sigma$ \emph{path} test, with the MED plug-in, fails to reject
\emph{any} pair ($0.92$, $0.51$, $0.81$) --- conditioning too restrictively, and
with an inflated $\hat\sigma$. The known-$\sigma$ \emph{union} recovers the
clearest split ($2$ vs $3$: $0.81 \to 1.6\times10^{-14}$) while staying cautious
on the closer pairs ($0.063$, $0.27$). The studentized test, needing no $\sigma$,
rejects all three ($2.3\times10^{-5}$, $0.028$, $1.4\times10^{-5}$).

\begin{table}[t]
\centering
\caption{Palmer penguins, $K=3$, $p$-values for the three cluster pairs
(\texttt{penguins\_pvalues\_k3.csv}). The naive test is absurd; the path test
rejects none; the union recovers pair $2$ vs $3$; the studentized test rejects
all three. Figure~\ref{fig:penguins}.}
\label{tab:penguins}
\begin{tabular}{lccccc}
\toprule
pair & naive & path (known) & union (known) & studentized $F$ \\
\midrule
1 vs 2 & $3.3\times10^{-57}$  & 0.920 & 0.063              & $2.3\times10^{-5}$ \\
1 vs 3 & $7.5\times10^{-29}$  & 0.514 & 0.271              & 0.028 \\
2 vs 3 & $5.6\times10^{-134}$ & 0.808 & $1.6\times10^{-14}$ & $1.4\times10^{-5}$ \\
\bottomrule
\end{tabular}
\end{table}

\subsection{Single-cell RNA sequencing (Pollen)}
We cluster $216$ developing-cortex cells from \citet{pollen2014lowcoverage}
(balanced across four inferred cell types; CP10k normalization, log$1$p, top-500
highly variable genes, ten standardized PCs) into $K=3$ groups and test the three
pairs (\texttt{scrna\_pvalues\_k3.csv}). Here the union is the hero: the original
\emph{path} test rejects \emph{none} of the three pairs ($0.475$, $0.179$,
$0.067$), whereas the \emph{union} rejects \emph{all three} ($0.0061$, $0.041$,
$0.043$). The studentized test rejects the clearest pair ($1$ vs $2$: $0.024$),
and the naive $p$-values are again absurd ($10^{-34}$, $10^{-31}$, $10^{-27}$).
Penguins and scRNA are complementary: on penguins the studentized test rejects
the most pairs, on scRNA the known-$\sigma$ union does --- but in both cases the
original path test is too conservative to be useful.

\begin{table}[t]
\centering
\caption{Pollen scRNA-seq, $K=3$, $p$-values for the three cluster pairs
(\texttt{scrna\_pvalues\_k3.csv}). The path test rejects none; the union rejects
all three. Figure~\ref{fig:scrna}.}
\label{tab:scrna}
\begin{tabular}{lcccc}
\toprule
pair & naive & path (known) & union (known) & studentized $F$ \\
\midrule
1 vs 2 & $4.6\times10^{-34}$ & 0.475 & 0.0061 & 0.024 \\
1 vs 3 & $3.5\times10^{-31}$ & 0.179 & 0.041  & 1.000 \\
2 vs 3 & $1.6\times10^{-27}$ & 0.067 & 0.043  & 1.000 \\
\bottomrule
\end{tabular}
\end{table}

% ===========================================================================
\section{Discussion}
\label{sec:discussion}
% ===========================================================================

We have given a more powerful selective test for $k$-means clustering that
conditions on the two tested clusters rather than the full Lloyd path, an exact
studentized version that needs no estimate of $\sigma$ and extends
\citet{yun2023unknownvar} to any $K$, and a fast solver for both. The
$\sigma$-free property is the practical headline: the most common failure of
plug-in approaches --- including data thinning's split step --- is that a global
variance estimate is inflated by exactly the cluster separation one is trying to
detect, and the studentized test sidesteps this entirely.

Several limitations remain. The studentized test requires the pointwise null
$H_0'$ (equal means \emph{within} each tested cluster), stronger than the equal-
cluster-means null; this is the inherent cost of not knowing $\sigma$, and a
violation is conservative unless within-cluster heterogeneity deflates the
residual under a true mean difference. The union's power gain is a known-$\sigma$
phenomenon and is ``studentized away'' under the $F$ pivot, because $k$-means
minimizes the within-scatter and pins the statistic against its floor. The method
is specific to $k$-means / Lloyd's algorithm, whereas \citet{yun2023unknownvar}
and \citet{gao2024hierarchical} cover broader clustering families.
% TODO: future work -- (i) a valid more-powerful union that actually gains under
% unknown variance (would require leverage that survives studentization);
% (ii) multiple-pairs / FWER control across cluster pairs, building on
% Yun & He (2024); (iii) randomized / data-fission variants for extra power
% (Rasines & Young 2023; Dharamshi et al. 2024).

% ===========================================================================
\section*{Figures}
% Figures are referenced by filename; PNGs live in ../sims/results/.
% Compile-time note: each \includegraphics points at a real PNG in the repo.
% ===========================================================================

\begin{figure}[p]\centering
  \includegraphics[width=0.85\textwidth]{decomposition.png}
  \caption{The contrast plane $(\Fnorm{\Pp_0\X},\Fnorm{\Pp_1\X})$. The
  known-$\sigma$ $\phi$-ray (fix within-scatter, vary separation $\to\chi$) versus
  the studentized $R$-fiber arc (fix total energy $T$, vary the angle $\to F$).
  Source: \texttt{sims/plot\_decomposition.R} $\to$ \texttt{decomposition.png}.}
  \label{fig:decomposition}
\end{figure}

\begin{figure}[p]\centering
  \includegraphics[width=0.85\textwidth]{cw_power.png}
  \caption{Known-$\sigma$ power: union (dashed) vs.\ path (solid) against
  separation, faceted by $q$. The union is uniformly more powerful
  (Table~\ref{tab:power}); the gain widens with $K$.
  Source: \texttt{cw\_power.png}.}
  \label{fig:cwpower}
\end{figure}

\begin{figure}[p]\centering
  \includegraphics[width=0.85\textwidth]{repro_yunbarber.png}
  \caption{Reproducing the Yun--Barber regime on the same truncation set:
  oracle $\ge$ studentized $F$ $\ge$ plug-ins. The studentized $F$ beats the
  sample plug-in once clusters separate, because the global SD absorbs the
  separation. Source: \texttt{repro\_yunbarber.png}.}
  \label{fig:yunbarber}
\end{figure}

\begin{figure}[p]\centering
  \includegraphics[width=0.49\textwidth]{cw_typeI.png}\hfill
  \includegraphics[width=0.49\textwidth]{gencov_typeI.png}
  \caption{Type I control. Left: QQ plot, selective tests on the diagonal, naive
  invalid (\texttt{cw\_typeI.png}). Right: under AR$(0.5)$/AR$(0.9)$ covariance,
  the general-covariance union/path stay calibrated (\texttt{gencov\_typeI.png}).}
  \label{fig:gencov}
\end{figure}

\begin{figure}[p]\centering
  \includegraphics[width=0.85\textwidth]{heavytail_typeI.png}
  \caption{Type I under heavy-tailed $t_5$/$t_{10}$ noise. The studentized test is
  fairly robust; mild anti-conservatism only for the heaviest tails.
  Source: \texttt{heavytail\_typeI.png}.}
  \label{fig:heavytail}
\end{figure}

\begin{figure}[p]\centering
  \includegraphics[width=0.85\textwidth]{penguins.png}
  \caption{Palmer penguins, $K=3$ (Table~\ref{tab:penguins}). Source:
  \texttt{penguins.png}.}
  \label{fig:penguins}
\end{figure}

\begin{figure}[p]\centering
  \includegraphics[width=0.85\textwidth]{scrna_k3.png}
  \caption{Pollen scRNA-seq, $K=3$ (Table~\ref{tab:scrna}): the union rejects all
  three pairs, the path rejects none. Source: \texttt{scrna\_k3.png}.}
  \label{fig:scrna}
\end{figure}

% ===========================================================================
\bibliographystyle{plainnat}
\bibliography{refs}
% ===========================================================================

\end{document}
